\chapter{Finite volume method (FVM)}
\label{chap:finite-volume-method}

\begin{quote}
Spec dir:
\passthrough{\lstinline!specs/2026-05-04-phase-14-finite-volume-method/!}.
Version: \passthrough{\lstinline!0.14.0!}.
\end{quote}

\subsection{What this gives you}\label{phase14-what}

A conservative cell-flux discretization of the diffusion operator with
\textbf{heterogeneous diffusivity} \(D(\mathbf x)\) as the headline new
capability:
\[
\partial_t\psi \;=\; \nabla\cdot\!\big(D(\mathbf x)\,\nabla\psi\big).
\]
The Phase 1 / Phase 5 / Phase 6 surface (FD
\passthrough{\lstinline!diff::laplacian!} +
constant-\(D\) \passthrough{\lstinline!solve::diffuse(\(\psi\), D, dt, n, tag)!})
is unchanged --- callers with constant \(D\) see no change. The new
heterogeneous-\(D\) path is the only place where FVM produces
different physics from FD; on uniform Cartesian grids the two paths
agree to round-off (locked in by the
\passthrough{\lstinline!AgreesWithFdLaplacianAtConstantD!} test gate).

Three new public symbols:

\begin{itemize}
\tightlist
\item
  \passthrough{\lstinline!gridcalc::fvm::cellLaplacian(psi, D)!} ---
  returns
  \(\nabla\cdot(D\,\nabla\psi)\) via cell-flux discretization with
  face-centered harmonic-mean \(D\)-averaging. New header
  \passthrough{\lstinline!include/gridcalc/fvm/CellLaplacian.hpp!}.
\item
  \passthrough{\lstinline!solve::diffuse(psi, D\_field, dt, n, tag)!}
  --- new heterogeneous-\(D\) overload. Reuses the existing
  \passthrough{\lstinline!solve::ExplicitEuler\{\}!} /
  \passthrough{\lstinline!solve::RK4\{\}!} integrator tags. CFL bound
  generalizes to
  \(D_{\max}\cdot dt\cdot\sum_a (1/h_a^2)\le \texttt{Tag::diffusionCFLLimit}\).
\item
  \textbf{Textbook demo.}
  \passthrough{\lstinline!examples/heterogeneous\_diffusion.cpp!} ---
  runnable CLI demo solving heterogeneous-\(D\) diffusion of a Gaussian
  peak through a smooth \(D(x) = D_0 + \varepsilon\cos(2\pi x/L)\)
  field. Produces VTK snapshots; per-snapshot diagnostics printed to
  stdout.
\end{itemize}

\subsection{When FVM, when FD?}\label{phase14-when}

Both paths discretize the Laplacian to second order on a uniform
Cartesian grid; for a constant \(D\) on a periodic domain they produce
the \emph{same} stencil. The choice matters only when
\(D\) varies in space (or when non-periodic boundary conditions are
in play; Phase 19 will land Dirichlet / Neumann walls and consume the
FVM foundation).

\begin{itemize}
\tightlist
\item
  Use \textbf{FD}
  (\passthrough{\lstinline!diff::laplacian!},
  constant-\(D\)
  \passthrough{\lstinline!solve::diffuse!}) for differential operators
  on uniform fields and constant-coefficient diffusion. Lower
  per-step cost; matches every test from Phases 1--6.
\item
  Use \textbf{FVM}
  (\passthrough{\lstinline!fvm::cellLaplacian!},
  heterogeneous-\(D\)
  \passthrough{\lstinline!solve::diffuse!}) for diffusion problems
  with a spatially varying \(D(\mathbf x)\). The cell-flux formulation
  is \emph{conservative by construction} and recovers the correct
  steady-state flux across \(D\)-discontinuities via the harmonic-mean
  face averaging.
\end{itemize}

\subsection{Cell-flux discretization}\label{phase14-cellflux}

At each cell \((i, j, k)\), the operator computes the net flux out
through the six faces, divided by the cell volume. For axis \(a\):
\[
[\nabla\!\cdot\!(D\nabla\psi)]_{a,(i,j,k)}
\;=\;
\frac{1}{h_a^2}\Big[\,D_{a,+}\,\big(\psi_{a,+}-\psi\big)
\;-\; D_{a,-}\,\big(\psi-\psi_{a,-}\big)\,\Big],
\]
with face-centered diffusivity computed as the \textbf{harmonic mean}
of the two adjacent cell values:
\[
D_{a,+} \;=\; \frac{2\,D(i)\,D(i+1)}{D(i) + D(i+1)},
\qquad
D_{a,-} \;=\; \frac{2\,D(i-1)\,D(i)}{D(i-1) + D(i)}.
\]
The harmonic mean is the standard FVM choice for diffusion (LeVeque
2007, Patankar 1980): it recovers the correct steady-state flux when
\(D\) jumps across a face, where the arithmetic mean would bias
toward the higher-\(D\) side. The implementation requires
\(D > 0\) at every cell; the harmonic mean is undefined otherwise.
The \passthrough{\lstinline!solve::diffuse!} heterogeneous-\(D\) overload
performs this check up-front and throws on violation.

\subsection{API}\label{phase14-api}

The Phase 1 / Phase 5 surface is unchanged. Phase 14 adds:

\begin{lstlisting}[language={C++}]
namespace gridcalc::fvm {

// $\nabla\cdot(D \nabla\psi)$ via cell-flux discretization with
// harmonic-mean face-D averaging. Periodic BCs only; D > 0 required.
core::Field<double> cellLaplacian(const core::Field<double>& psi,
                                  const core::Field<double>& D);

}  // namespace gridcalc::fvm

namespace gridcalc::solve {

// New heterogeneous-D overload. Reuses ExplicitEuler / RK4 tags.
template <class Tag = ExplicitEuler>
void diffuse(core::Field<double>& psi,
             const core::Field<double>& D,
             double dt,
             int n_steps,
             Tag tag = {});

}  // namespace gridcalc::solve
\end{lstlisting}

\subsection{Demo: Gaussian peak in a heterogeneous
medium}\label{phase14-demo}

Build with \passthrough{\lstinline!GRIDCALC\_BUILD\_EXAMPLES=ON!} (the
same flag Phase 12 added):

\begin{lstlisting}[language=bash]
cmake --preset clang-release -DGRIDCALC_BUILD_EXAMPLES=ON
cmake --build --preset clang-release --target heterogeneous_diffusion
./build/clang-release/examples/heterogeneous_diffusion \
    --n-x 64 --n-steps 4000 --snapshot-every 200 \
    --d-mean 0.10 --d-amplitude 0.07 \
    --gaussian-sigma 0.5 --integrator rk4 --out-dir ./out
\end{lstlisting}

This is the run that produced the figures below. The Gaussian peak is
centered on a 64\(^3\) cubic box; \(D(x) = 0.10 + 0.07\cos(2\pi x/L)\)
varies smoothly along \(x\), with maxima at \(x \in \{0, L\}\) and a
minimum at \(x = L/2\). The Gaussian spreads anisotropically: faster
in the high-\(D\) regions, slower in the low-\(D\) region. Per-snapshot
diagnostics on stdout track mass conservation, \(L^2\)-norm decay,
positivity, and peak amplitude.

\begin{figure}[h]
  \centering
  \includegraphics[width=0.32\textwidth]{figures/heterogeneous-diffusion/early.png}\hfill
  \includegraphics[width=0.32\textwidth]{figures/heterogeneous-diffusion/mid.png}\hfill
  \includegraphics[width=0.32\textwidth]{figures/heterogeneous-diffusion/late.png}
  \caption{$z$-midplane slices of $\psi$ from a 64\(^3\)
  heterogeneous-$D$ diffusion run with $D(x) = 0.10 + 0.07\cos(2\pi
  x/L)$, Gaussian peak IC ($\sigma = 0.5$, centered),
  $\Delta t \approx 0.024$, RK4. The peak preferentially diffuses
  into the high-$D$ corridor along $x$ while staying narrow along
  the low-$D$ axis at the center; peak amplitude monotonically
  decays as the Gaussian relaxes toward the steady state $\psi = \langle\psi_0\rangle$.}
  \label{fig:heterogeneous-diffusion}
\end{figure}

The Phase 14 acceptance test
(\href{../../../test/heterogeneous_diffusion_test.cpp}{\passthrough{\lstinline!test/heterogeneous\_diffusion\_test.cpp!}})
runs a stripped-down version of the same simulation in CI and asserts
mass conservation, \(L^2\)-norm monotone decrease, and positivity
preservation. The
\passthrough{\lstinline!fvm::cellLaplacian!} unit tests
(\href{../../../test/fvm_cell_laplacian_test.cpp}{\passthrough{\lstinline!test/fvm\_cell\_laplacian\_test.cpp!}})
verify the harmonic-mean face flux on a 4-cell jump test, the
order-2 convergence on a manufactured \(D(x) = 1 + 0.5\cos(x)\),
\(\psi(x) = \sin(x)\) pair, and the conservative property
\(\sum \texttt{cellLaplacian}(\psi, D) = 0\) on the periodic domain.

\subsection{What this does \emph{not} do}\label{phase14-not-do}

\begin{itemize}
\tightlist
\item
  \textbf{No FD replacement.} Phase 1's
  \passthrough{\lstinline!diff::laplacian!}, Phase 2's
  \passthrough{\lstinline!gradient!} /
  \passthrough{\lstinline!divergence!}, and Phase 5 / Phase 6's
  constant-\(D\)
  \passthrough{\lstinline!solve::diffuse!} are unchanged.
\item
  \textbf{No higher-order FVM.} 2nd-order central-difference face
  fluxes only; MUSCL / WENO reconstructions are out of scope for the
  v1.0 roadmap.
\item
  \textbf{No vector / tensor FVM.} Phase 13's vector / tensor
  \passthrough{\lstinline!Field!} surface stays separate.
\item
  \textbf{No non-periodic BCs.} Phase 19 (the renumbered
  Dirichlet / Neumann phase) handles them and consumes the FVM
  foundation as the natural starting point.
\item
  \textbf{No implicit / IMEX time stepping.} Phase 20 (renumbered)
  covers implicit schemes; the heterogeneous-\(D\) overload accepts
  the explicit \passthrough{\lstinline!ExplicitEuler\{\}!} /
  \passthrough{\lstinline!RK4\{\}!} tags only.
\end{itemize}

For the harmonic-mean derivation, the order-2 truncation analysis,
the generalized CFL bound, and the rationale for keeping the FVM /
FD paths separate, see the ``Finite volume method'' chapter of the
Developer Note.
